\section{Appendix}

    \subsection*{Derivation of Malus's Law}
    \label{sec:derivationmalus}

Consider linearly polarized light with intensity $I_0$ and polarization direction $\vec{e}_P$
incident on an analyzer whose transmission axis $\vec{e}_A$ is rotated by an angle $\theta$
relative to $\vec{e}_P$.
Only the component of the electric field along $\vec{e}_A$ is transmitted.

The incident electric field can be written as
\begin{equation}
    \vec{E} = E_0 \, \vec{e}_P,
\end{equation}
where $E_0$ is the amplitude. The transmitted component along $\vec{e}_A$ is
\begin{equation}
    E_t = \vec{E} \cdot \vec{e}_A = E_0 \cos\theta.
\end{equation}
Since the intensity is proportional to the square of the field amplitude, $I \propto E^2$, it follows that
\begin{equation}
    I(\theta) = I_0 \cos^2\!\theta
\end{equation}

\subsection*{Mean Intensity after the Analyzer}

If unpolarized light of intensity $I_\text{in}$ is incident on the polarizer,
an ideal polarizer transmits exactly half of it:
\begin{equation}
    I_0 = \frac{I_\text{in}}{2}.
\end{equation}
This follows from averaging $\cos^2\theta$ over all polarization directions:
\begin{equation}
    \langle \cos^2\theta \rangle
    = \frac{1}{2\pi} \int_0^{2\pi} \cos^2\!\theta \, \mathrm{d}\theta
    = \frac{1}{2\pi} \int_0^{2\pi} \frac{1 + \cos 2\theta}{2} \, \mathrm{d}\theta
    = \frac{1}{2}.
\end{equation}
Therefore, the intensity of the light incident on the analyzer is
\begin{equation}
    I_0 = \frac{I_\text{in}}{2}
\end{equation}
i.e.\ the intensity after the polarizer is exactly half of the incoming intensity,
independent of the orientation of the analyzer.